\section{Nonnegative merge deficits and spectral projections}
\label{sec:merge-deficit-positivity}

Let a finite group act by two commuting permutation representations on a
finite orthonormal basis, and let the combined action be their pointwise
product. Write $F_Q,F_E,F_{QE}$ for the logarithmic irreducible-dimension
observables of these actual actions. The dimension inequality in
\cite[Section~5, Lemma~6.1(3)]{OpenAI2026AreaLaw} gives the following
nonnegativity statement, used in the merge decomposition of Section~7.

\begin{theorem}[Nonnegative merge deficit]
\label{thm:merge-deficit-positivity}
\lean{PermutationRepresentation.posSemidef_mergeDeficit,
TensorPower.posSemidef_pairMergeDeficit}
\uses{thm:schur_label_merge_dimensions,def:pair_merge_deficit}
\leanok
The actual merge deficit is positive semidefinite:
\[
 F_Q+F_E-F_{QE}\ge0.
\]
In particular, this holds for the separate and simultaneous actions on the
literal paired copy space, at every number of copies.
\end{theorem}

\begin{proof}
\leanok
The three commuting families of central projections give a joint orthogonal
resolution. On the sector with labels $\lambda,\mu,\nu$, the deficit is
\[
 \log d_\lambda+\log d_\mu-\log d_\nu.
\]
Every nonzero joint sector satisfies $d_\nu\le d_\lambda d_\mu$, so its
coefficient is nonnegative. Vanishing sectors contribute zero. Positivity
therefore follows by summing the actual Hermitian projections with these
coefficients. This also covers a zero-dimensional ambient space.
\end{proof}

The physical expression $F_Q+F_V-F_Y$ has an independent middle region $Y$.
Its positivity on the physical symmetric subspace requires the complementary
label identity there. It is not a full-space instance of the preceding theorem.
The following elementary spectral fact applies once that identity has been
derived, as in \cite[Section~7, lines~515--522]{OpenAI2026AreaLaw}.

\begin{theorem}[Nonnegative spectral projection on an intertwined range]
\label{thm:nonnegative-spectral-intertwiner}
\lean{Matrix.PosSemidef.spectralProjectionGE_zero,
Matrix.spectralProjectionGE_zero_mul_of_intertwine}
\leanok
For a Hermitian matrix $A$, let $J_A=\mathbf1_{[0,\infty)}(A)$.
If $B\ge0$, then $J_B=I$. If $T$ is any rectangular matrix satisfying
$AT=TB$, then
\[
 J_AT=T.
\]
The matrix $T$ may have zero or noninjective range, and the coordinate sets
may be empty.
\end{theorem}

\begin{proof}
\leanok
All eigenvalues of $B$ are nonnegative, including its zero eigenvalues, so
the closed nonnegative spectral projection is the identity. Functional
calculus preserves the Hermitian intertwining relation, giving
$J_AT=TJ_B=T$.
\end{proof}
